\section{Threat Model and Security Goals}\label{sec:threat}

\subsection{Trust assumptions}
We assume that TLS is secure; that the bank server is honest but its database
may leak (read-only exposure); that the phone running the authenticator is
trusted; and that the environment is clean at the moment of device enrolment
(trust on first use). These assumptions are discussed, and their fragility noted,
in \cref{sec:discussion}.

\subsection{Attacker capabilities}
\Cref{tab:attackers} lists the attacker capabilities in scope. The scenarios in
the last column (\cref{sec:eval}) exercise each capability.

\begin{table}[htbp]
\centering
\caption{Attacker capabilities.}
\label{tab:attackers}
\small
\begin{tabularx}{\linewidth}{@{}clXl@{}}
\toprule
ID & Attacker & Capability & Scenarios \\
\midrule
A1 & Man-in-the-browser     & Reads and rewrites everything in the browser & S2a, S2b, S3, S4 \\
A2 & Real-time phishing proxy & Relays and rewrites traffic in real time & S2a, S2b, S3, S4 (see note) \\
A3 & Replay                 & Re-uses an old code or signature & S6, S7, S8 \\
A4 & Online guessing        & Guesses codes against the live server & S9 \\
A5 & Database theft         & Holds a copy of the server database & S10 \\
A6 & Malware on the phone   & Compromises the phone, or PC and phone together & out of scope \\
A7 & Social engineering     & Persuades the user to approve & out of scope \\
\bottomrule
\end{tabularx}
\end{table}

\paragraph{Note on A2 (equivalence to A1).}
We give the phishing proxy no separate scenario because, in this model, it can do
exactly what the MitB layer does: stand between the user and the bank, rewrite
the transfer request, and rewrite the returned page (including the QR code). The
only difference is how the attacker reaches that position (a fake site versus
malware), which does not change the outcome. Every claim about real-time
phishing in this report (RQ1) therefore refers to scenarios S2a/S2b/S3/S4 and
relies on this equivalence; we did \emph{not} run a separate phishing
experiment.

\paragraph{Note on A6 (phone compromise, out of scope).}
The assumption that the phone is trusted and separate from the PC is a strong
one, and the literature shows it is easily lost: cross-infection of PC and phone
\cite{dmitrienko2014}, cross-platform synchronisation that erases the gap between
the two devices \cite{konoth2016}, and the bank app and authenticator running on
the same phone \cite{haupert2018}. We treat A6 as out of scope and return to it
in \cref{sec:discussion}.

\subsection{Security goals}
\Cref{tab:goals} states the goals the modes are measured against. Two goals carry
limitations that are important enough to state up front.

\begin{table}[htbp]
\centering
\caption{Security goals and the scenarios that test them.}
\label{tab:goals}
\small
\begin{tabularx}{\linewidth}{@{}clX l@{}}
\toprule
ID & Goal & Meaning & Tested by \\
\midrule
G1 & Authorization integrity & Server executes $T$ only if the device produced a code for exactly $T$ & S4, S5 \\
G2 & Freshness               & No replay, no late use & S6, S7 \\
G3 & Single use              & One code, one execution, even under concurrency & S8 \\
G4 & Guessing limit          & Bounded guessing \emph{per transaction} (see limitation) & S9 \\
G5 & DB-leak resilience      & A leaked database does not let the attacker forge transfers & S10 \\
G6 & Verifiable evidence     & A third party can check the evidence (necessary, not sufficient, for non-repudiation) & S11 \\
G7 & WYSIWYS                 & The user saw exactly $T$ (idealised attentive user) & S2a, S3 \\
\bottomrule
\end{tabularx}
\end{table}

\paragraph{Limitation of G6 (verifiable evidence, not non-repudiation).}
Scenario S11 shows only that the evidence stored in the database can be checked
with a public key, without any secret of the bank's. It is \emph{not}
non-repudiation, for two reasons. First, the authenticator is a web application
served by the bank itself, so the bank could serve different JavaScript that uses
the signing key for an arbitrary transaction; a non-extractable key prevents
\emph{exporting} the key, not \emph{using} it. Second, the public key lives only
in the bank's \texttt{devices} table, with no independent party attesting that it
belongs to the user, so the bank could substitute a key. We therefore write
\emph{verifiable evidence}, not \emph{non-repudiation}, for S11, and mark the
\texttt{sig\_blind} cell with an asterisk (the signature proves the device
signed, not that the user saw the transaction).

\paragraph{Limitation of G4 (per-transaction, not per-user).}
The three-attempt limit is enforced per transaction (\texttt{server/verify.py},
\texttt{settings.max\_attempts}); there is no per-user or per-device counter, and
transaction creation is not rate-limited. An attacker with the password but no
phone (A4) can create $k$ transactions and guess three times each, so the
cumulative success probability is about $k\cdot p$. This limitation is quantified
in \cref{sec:eval}.
